\section{端云协同与端侧资源约束}
\label{sec:rt4}

本节作为第六部分，严格承接 T+8 第 4 章关于端云协同可行性的结论 C4-1、C4-2 与 C4-3，只讨论当前 Hermes Soul 与 ORCA Workforce 原型在 Demo 验证中新增暴露的问题，不重复论证端云协同本身是否成立。这里的增量重点有三个：端侧运行难点如何在真实多 Agent 流程中转化为资源约束，Hermes 接入后哪些问题推动了架构决策，以及这些决策如何与 RT2 的跨端恢复、RT6 的成本优化形成耦合。与第~\ref{sec:rt3}~节的分工一致，本节默认 Hermes 不是普通应用执行器，而是系统级 Soul 大脑；端侧执行仍由 Workforce 和专用 App Agent 承担。

专家反馈 F3 与 F4 要求补充的不是泛化的端云协同描述，而是多 Agent 架构在实际 Demo 中遇到的瓶颈、规避方案和可复现实验设计。基于当前代码和运行记录，本节将“管家问题”作为主要分析线索：当用户给出一句模糊指令时，系统入口是否能先理解、补全和提问，决定了后续执行层是稳定协同，还是把不完整任务继续下发给应用 Agent，造成空转、重复提问和资源浪费。

\subsection{端侧资源约束分析}

在 T+8 第 4 章结论 C4-1 至 C4-3 的基础上，本阶段的端侧资源约束需要从“端上算力有限”进一步细化为可影响架构选型的执行约束。端侧并非完全不能运行智能体，而是不适合长期承载多个重型通用大脑、长上下文推理和跨应用全量数据汇聚。对于 ORCA-Hermes 原型而言，端侧更适合保存应用数据、维护用户交互状态、执行确定性工具和呈现结果；云侧更适合承担复杂语义理解、任务富化、创作生成和高成本规划。这个分工不是把端侧降级为壳，而是在端侧保留数据边界和执行控制权，在云侧释放大模型推理能力。

当前原型的资源约束可以概括为表~\ref{tab:rt4-edge-constraints}。表中每一项都直接对应 Demo 中的工程决策：Hermes 被放在 Soul 层而不是复制到所有 App Agent，是为了控制多实例内存和上下文成本；App Agent 保留应用边界，是为了降低工具误调用和隐私扩散；统一运行环境与任务状态记录，则是为了让云侧 Soul 与端侧执行层稳定衔接。

\begin{table}[h]
\centering
\begin{tabular}{p{0.17\linewidth}p{0.36\linewidth}p{0.37\linewidth}}
\hline
\textbf{约束类型} & \textbf{对多 Agent Demo 的影响} & \textbf{当前架构中的处理方式} \\
\hline
算力与时延 & 端侧无法稳定承载多个长链路大模型推理，用户交互也不能等待无限规划 & Hermes 只承担 Soul 大脑，重推理上提到云侧，执行层保持轻量调度与工具调用 \\
内存与状态 & 多个重型实例会重复维护角色、记忆、会话和工具状态，恢复成本随实例数上升 & 不把 Hermes 扩散为所有 App Worker，执行层继续使用专用 App Agent \\
上下文规模 & 用户画像、应用观察、历史经历和工具结果叠加后容易形成高 Token 请求 & Soul 读取摘要和必要偏好，执行层仅回传与当前任务相关的结果 \\
工具边界 & 单 Agent 持有全量工具时容易误选工具，也难以定位失败责任 & Workforce 按应用边界分配任务，App Agent 只调用自身领域工具 \\
隐私与授权 & 相册、联系人、订单和笔记等个人数据不应无差别进入同一云侧上下文 & 端侧保留原始数据，云侧只接收任务所需摘要、偏好和最小必要结果 \\
网络与响应波动 & 云侧推理存在响应时间和可用性波动，可能影响入口交互体验 & 通过任务状态记录、进度反馈和确认机制降低等待感 \\
\hline
\end{tabular}
\caption{端侧多 Agent 执行的关键资源约束}
\label{tab:rt4-edge-constraints}
\end{table}

这些约束重新限定了“端侧化”的含义。端侧化不是把所有模型推理都搬回设备，也不是让一个单 Agent 在端上持有全部应用工具，而是让端侧保存用户数据、应用状态和执行确认，让云侧模型只在需要强理解和强生成时参与。由此形成的端云协同边界比单 Agent 模式更细：Soul 负责判断任务是否清楚、是否需要用户授权、是否应补充画像上下文；Workforce 负责把清楚的任务拆解给具体 App Agent；App Agent 负责在自己的应用边界内执行。

\subsection{难点与解决思路}

Demo 测试中最典型的问题是“管家式入口”没有足够早地拦截模糊任务。用户输入“帮我发一个小红书”时，这句话本身缺少主题、素材、语气、是否带图和是否允许发布等关键约束。如果 Soul 只是把原话或一段澄清文本交给 Workforce，执行层会继续创建小红书任务，再由小红书 Agent 承担追问职责。此时系统并不是工具不可用，而是协同位置不合理：用户面对的是一个系统管家，澄清应该发生在大脑层，而不是发生在已经进入应用执行后的 Worker 层。

该问题推动了当前原型的核心架构选择。Hermes Soul 负责在入口处理解和丰富任务，识别外部可见动作、创作型任务和持久化修改是否需要用户确认；Workforce 只接收可执行的 enriched task，而不是接收原始模糊请求。这样做可以减少执行层空转，也可以把用户确认和授权记录放在任务开始前，避免 App Agent 在缺少上下文时自行猜测。表~\ref{tab:rt4-demo-painpoints} 汇总了 Demo 过程中暴露出的具体问题和相应处理方向。

\begin{table}[h]
\centering
\begin{tabular}{p{0.22\linewidth}p{0.33\linewidth}p{0.33\linewidth}}
\hline
\textbf{测试中暴露的问题} & \textbf{对体验与资源的影响} & \textbf{架构处理方向} \\
\hline
模糊发布请求过早下发 & 执行层被迫追问，任务链路变长，用户感知为系统没有真正理解意图 & 将澄清、补全和授权判断前移到 Hermes Soul \\
上下文过大 & 用户画像、历史经历和应用观察若不加裁剪，会放大 Token 成本和入口时延 & 使用画像摘要、按需检索和任务结果摘要控制上下文规模 \\
端侧工具分散 & 相册、联系人、订单、笔记和发布能力分布在不同 App 中，单入口难以直接执行全部动作 & 由 Workforce 将 enriched task 拆给具备对应工具边界的 App Agent \\
多应用状态衔接 & 跨应用任务需要记录每一步是否完成、是否确认和是否需要恢复 & 通过任务状态记录支持中断恢复、结果汇总和后续追踪 \\
用户等待时间 & 入口理解、用户澄清和应用执行串联后，用户可能感知为响应变慢 & 在 Soul 层合并关键问题，并通过执行进度反馈降低等待感 \\
外部可见动作缺少确认 & 发布、发消息等动作若直接执行，会带来用户信任和安全风险 & 在 Soul 层识别高影响动作，在执行前保留必要确认 \\
\hline
\end{tabular}
\caption{Demo 验证中的问题、影响与架构处理方向}
\label{tab:rt4-demo-painpoints}
\end{table}

解决思路需要围绕“入口理解”和“执行边界”展开，而不是单纯增加更多 Agent。复杂理解、任务富化、跨应用意图判断和高风险动作识别由 Hermes Soul 承担；确定性工具调用、应用数据读写和结果沉淀由端侧 App Agent 承担。对于照片理解、旅行内容整理和小红书创作这类高成本任务，系统应先在端侧缩小数据范围，再将必要摘要交给云侧模型生成内容。对于联系人消息、发布内容和画像修改这类外部可见或持久化动作，系统应在 Soul 层形成确认语义，再由执行层完成具体操作。

Token 消耗优化也应嵌入这套流程。Soul 不应把完整画像、全部经历和所有应用数据一次性塞入模型，而应读取稳定摘要、近期相关经历和当前任务所需的 App 偏好观察。Workforce 不应把每个中间结果原样回灌给大脑，而应将工具结果压缩成可恢复、可评估、可继续执行的任务摘要。这样，RT6 中的成本优化就不只是更换模型或降低调用频率，而是通过上下文裁剪、缓存复用、失败重试控制和工具边界减少无效 Token。

\subsection{端云协同作为解法：Hermes 运行效果实证}

本阶段能够复用的实证数据来自 Hermes Soul 的运行记录和 Workforce 的任务事件记录。相同模糊任务“帮我发一个小红书”中，Hermes Soul 能够识别任务缺少主题与内容方向，并返回需要用户补充的信息；对应运行记录显示 Hermes Soul 用时约 23.4 秒，事件数为 5，事件序列包含命令可用性、思考片段、文本输出和完成信号。这个结果说明 Hermes 作为云侧 Soul 大脑能够发现任务缺口，但也暴露出入口澄清的时延仍然偏高，后续需要通过更短的系统上下文、更明确的澄清策略和更直接的人机交互降低首轮等待。

表~\ref{tab:rt4-hermes-observation} 给出当前可写入报告的观测值。这里需要区分“已观测结果”和“后续对照计划”：Hermes 的运行效果已有日志证据，OpenClaw 与 ReAct+Skill 单 Agent 的同案对照还需要在相同任务集下补测，因此本节只给出对照协议，不把尚未完成的数据写成结论。

\begin{table}[h]
\centering
\begin{tabular}{p{0.23\linewidth}p{0.23\linewidth}p{0.42\linewidth}}
\hline
\textbf{观测项} & \textbf{记录值} & \textbf{解释} \\
\hline
模糊小红书发布请求 & Hermes Soul 约 23.4 秒返回澄清内容，事件数为 5 & 云侧 Soul 能识别任务缺口，但入口时延和交互阻塞策略仍需优化 \\
事件流结构 & 包含可用命令、思考片段、文本输出和完成信号 & Hermes 具备过程可观测性，后续可与 Workforce 日志合并形成端到端任务记录 \\
MCP 工具桥注册 & Soul 工具桥注册用户画像、任务状态、画像更新、经历追加和用户提问能力 & Hermes 能接入端侧画像与交互，但工具集合需要保持最小必要 \\
执行层任务流 & 模糊任务被创建并分配给小红书相关 Agent & Workforce 能继续调度 App Agent，但澄清不应滞留到执行层 \\
\hline
\end{tabular}
\caption{Hermes Soul 端云协同运行观测}
\label{tab:rt4-hermes-observation}
\end{table}

这些数据支撑的结论比较明确：Hermes 接入后，系统具备了在入口处识别模糊任务的能力；但如果澄清只以普通文本返回，且没有阻塞式吸收用户回复，任务仍可能被执行层误解为待执行内容。因而端云协同的关键不只是“云侧模型是否足够强”，还包括“云侧 Soul 的判断如何影响端侧任务状态”。更合理的闭环是 Hermes Soul 发现高影响缺口后直接通过用户交互通道提问，待用户补充后再形成 enriched task；Workforce 只处理已经满足执行条件的任务。

为回应 F3/F4 对竞品对比的要求，后续对照测试应采用同一任务、同一模拟数据、同一确认规则和同一评价指标。表~\ref{tab:rt4-compare-protocol} 给出建议协议，用于后续补充视频和报告数据。这里的 ORCA-Hermes 多 Agent 指当前 Hermes Soul 指挥 Camel Workforce 的原型形态，OpenClaw 与 ReAct+Skill 单 Agent 作为对照组，用于验证多 Agent 架构在跨应用任务、工具边界和失败恢复上的差异。

\begin{table}[h]
\centering
\begin{tabular}{p{0.18\linewidth}p{0.25\linewidth}p{0.25\linewidth}p{0.20\linewidth}}
\hline
\textbf{同案任务} & \textbf{ORCA-Hermes 多 Agent 观察点} & \textbf{OpenClaw / ReAct+Skill 单 Agent 观察点} & \textbf{记录指标} \\
\hline
模糊发布任务 & Soul 是否先澄清主题、素材、授权，再交给小红书 Agent & 单 Agent 是否能自行澄清并避免误发布 & 首次澄清时延、总轮数、误调用次数 \\
旅行内容发布 & Soul 是否把相册、旅行信息和小红书创作拆给对应 App Agent & 单 Agent 是否在全量工具中稳定完成跨应用链路 & Token 总量、工具调用数、完成率 \\
联系人消息任务 & Soul 是否识别外部可见动作并保留发送确认 & 单 Agent 是否同时正确处理联系人查找和发送风险 & 确认次数、误发送风险、恢复能力 \\
资源受限场景 & 端侧是否保留数据与执行，云侧只承担理解和生成 & 单 Agent 是否因长上下文和全量工具导致成本或失败上升 & 峰值上下文、端侧等待时间、失败原因 \\
\hline
\end{tabular}
\caption{相同任务下的多 Agent 与单 Agent 对照测试协议}
\label{tab:rt4-compare-protocol}
\end{table}

\begin{figure}[h]
\centering
\fbox{
\begin{minipage}{0.88\linewidth}
\centering
\textbf{端云协同验证闭环}\\[0.5em]
端侧用户请求
$\rightarrow$ Hermes Soul 判断任务缺口
$\rightarrow$ 用户澄清或授权
$\rightarrow$ enriched task
$\rightarrow$ Workforce 调度
$\rightarrow$ App Agent 执行
$\rightarrow$ 结果记录与成本统计
\end{minipage}
}
\caption{Hermes Soul 端云协同验证闭环}
\label{fig:rt4-edge-cloud-flow}
\end{figure}

\subsection{端侧落地的关键约束与对立项的资源含义}

端侧落地的关键约束集中在几组资源对立关系上。Hermes 提供更强的长期记忆、人格化上下文和任务理解能力，但它本身比普通 App Agent 更重，因此不能简单复制为多个端侧执行器。当前原型选择让 Hermes 只作为 Soul 大脑，是对强认知能力与低资源约束之间的折中：系统保留一个强入口，避免每个应用内部都维护一套重型认知状态。

长上下文与低成本之间也存在直接冲突。用户画像、经历、应用偏好和工具结果越完整，模型越容易理解用户真实意图，但 Token、时延和隐私风险也会随之上升。RT6 的成本优化需要沿着这条约束继续推进，把全量上下文改为按需摘要，把每轮工具结果改为任务级摘要，把失败重试限制在必要范围内。只有在上下文被显式治理后，端云协同才会真正降低端侧压力，而不是把端侧问题转化为云侧调用成本。

端侧隐私与云侧推理能力之间的关系决定了数据流向。云侧 Hermes 适合处理模糊意图、创作生成和跨应用语义判断，但相册、联系人、订单和笔记等原始数据不应无差别上送。当前架构中 App Agent 继续持有应用边界，Soul 只读取摘要、偏好和必要观察，是为了让云侧模型获得足够语义信息，同时不破坏端侧数据隔离。这个边界也使后续跨端恢复更容易描述：RT2 只需要恢复任务状态、用户确认状态和必要摘要，而不是迁移所有应用原始数据。

用户体验与安全确认之间的矛盾最终落在任务状态设计上。发布小红书、发送消息、修改长期画像等动作需要确认，但过多确认会破坏管家式体验。当前原型的方向是把高影响澄清前移到 Hermes Soul，把具体执行确认保留在 App Agent 之前，使系统既能主动理解和补全任务，也能在外部可见动作前保留用户控制权。当任务跨设备迁移或中断恢复时，系统必须明确记录“待澄清、已澄清、待确认、执行中、已完成”等状态，否则很容易出现重复提问、重复执行或漏确认的问题。

综合来看，第六部分要强调的架构决策并不是“端侧无法运行智能体，所以全部交给云端”，而是“端侧保留数据、工具和确认，云侧承担强理解和生成，Workforce 负责把二者连接成可恢复的任务流”。Hermes 的价值在于作为 Soul 大脑提前发现任务缺口和资源风险，Camel Workforce 与 App Agent 的价值在于保持应用执行边界。该设计同时回应了 F3 关于 Token 消耗和端侧难点的要求，也回应了 F4 关于 Demo 瓶颈、架构选型依据和竞品对照测试设计的要求。
